%% chapter5.tex – Design and Implementation
%% Group No. 2 | Dept. of CSE (Data Science), MACE Kothamangalam
%% Team: Muhammed Ajmal P (MAC23CD040), Muhammed Adil PP (MAC23CD039),
%%       Mohammed Shadeed P (MAC23CD037), Sinan Rahman (LMAC23CD066)
%% Guide: Prof. Eldo P Elias | May 2026

\chapter{Design and Implementation}
\label{chap:design}

\section{System Architecture}

WaveGuard~V4 follows a three-layer architecture as shown in
Figure~\ref{fig:sysarch}:

\begin{figure}[H]
\centering
\begin{tikzpicture}[
  node distance=1.4cm and 2.2cm,
  box/.style={draw, rounded corners, minimum width=3.8cm, minimum height=0.9cm,
              align=center, font=\small},
  arr/.style={-Stealth, thick},
  layer/.style={draw=gray!60, fill=gray!8, rounded corners, inner sep=8pt}
]
% Layer 1 – Sensors
\node[box, fill=blue!15]  (mpu)  {MPU6050\\(50 Hz accel)};
\node[box, fill=blue!15, right=1.8cm of mpu] (gps) {NEO-6M\\(GPS 1 Hz)};

% Layer 2 – MCU
\node[box, fill=orange!20, below=1.4cm of mpu, xshift=2.2cm]
      (esp) {ESP32\\DevKit v1\\(240 MHz)};

% Layer 3 – Output
\node[box, fill=green!20, below=1.4cm of esp, xshift=-2.2cm]
      (wifi) {WiFi\\Dashboard\\(React)};
\node[box, fill=green!20, right=1.8cm of wifi]
      (gsm)  {SIM800L\\SMS Alert};
\node[box, fill=green!20, right=1.8cm of gsm]
      (db)   {FastAPI\\+ SQLite};

% Arrows
\draw[arr] (mpu) -- node[left, font=\scriptsize]{I2C} (esp);
\draw[arr] (gps) -- node[right, font=\scriptsize]{UART} (esp);
\draw[arr] (esp) -- node[left, font=\scriptsize]{HTTP} (wifi);
\draw[arr] (esp) -- node[font=\scriptsize, above]{AT} (gsm);
\draw[arr] (esp) -- node[right, font=\scriptsize]{REST} (db);

% Layer labels
\node[font=\bfseries\small, above=0.2cm of mpu, xshift=2cm]{Layer 1: Hardware Sensors};
\node[font=\bfseries\small, below=0.1cm of esp]{Layer 2: ESP32 Processing};
\node[font=\bfseries\small, below=0.2cm of wifi, xshift=2.8cm]{Layer 3: Output};
\end{tikzpicture}
\caption{WaveGuard V4 System Architecture (three-layer IoT model)}
\label{fig:sysarch}
\end{figure}

\section{Pin Configuration}

Table~\ref{tab:pins} lists the complete GPIO assignment for the ESP32 DevKit~v1.

\begin{table}[H]
\centering
\caption{ESP32 GPIO Pin Assignment}
\label{tab:pins}
\begin{tabular}{llll}
\toprule
\textbf{GPIO} & \textbf{Function}   & \textbf{Connected To}   & \textbf{Protocol} \\
\midrule
GPIO21        & SDA                 & MPU6050 pin 4 (SDA)     & I2C (4.7~k$\Omega$ pull-up) \\
GPIO22        & SCL                 & MPU6050 pin 3 (SCL)     & I2C (4.7~k$\Omega$ pull-up) \\
GPIO27        & UART TX             & SIM800L RX              & UART 9600 \\
GPIO26        & UART RX             & SIM800L TX              & UART 9600 \\
GPIO16        & UART2 RX            & NEO-6M TX               & UART 9600 \\
GPIO17        & UART2 TX            & (not connected)         & UART 9600 \\
3V3           & 3.3 V supply        & MPU6050 VCC, NEO-6M 3V3 & Power \\
5V (Vin)      & 5 V supply          & MT3608 output           & Power \\
GND           & Ground              & All modules             & Common GND \\
\bottomrule
\end{tabular}
\end{table}

\section{Data Flow Diagram}

Figure~\ref{fig:dfd} illustrates the level-0 data flow from sensor input to alert
output.

\begin{figure}[H]
\centering
\begin{tikzpicture}[
  node distance=0.9cm and 2.0cm,
  proc/.style={draw, ellipse, minimum width=2.8cm, minimum height=0.9cm,
               align=center, font=\small, fill=yellow!25},
  stor/.style={draw, rectangle, minimum width=2.8cm, minimum height=0.7cm,
               align=center, font=\small, fill=cyan!20},
  ext/.style ={draw, rectangle, minimum width=2.2cm, minimum height=0.7cm,
               align=center, font=\small, fill=gray!20},
  arr/.style ={-Stealth, thick}
]
\node[ext]  (sensor) {MPU6050};
\node[proc, right=2cm of sensor] (sample) {Sample\\$a_x,a_y,a_z$};
\node[proc, right=2cm of sample] (rms)    {Compute\\RMS};
\node[proc, right=2cm of rms]    (class)  {Classify\\State};
\node[ext,  below=1.5cm of class] (sms)   {SIM800L\\SMS};
\node[stor, below=1.5cm of rms]   (db)    {SQLite\\DB};
\node[ext,  below=1.5cm of sample](dash)  {React\\Dashboard};

\draw[arr] (sensor) -- node[above,font=\scriptsize]{I2C} (sample);
\draw[arr] (sample) -- node[above,font=\scriptsize]{50 Hz} (rms);
\draw[arr] (rms)    -- node[above,font=\scriptsize]{RMS} (class);
\draw[arr] (class)  -- node[right,font=\scriptsize]{HIGH\_WAVE} (sms);
\draw[arr] (rms)    -- node[right,font=\scriptsize]{POST} (db);
\draw[arr] (db)     -- node[below,font=\scriptsize]{GET} (dash);
\end{tikzpicture}
\caption{WaveGuard V4 Level-0 Data Flow Diagram}
\label{fig:dfd}
\end{figure}

\section{Power System Design}

Figure~\ref{fig:power} shows the power distribution tree.

\begin{figure}[H]
\centering
\begin{tikzpicture}[
  box/.style={draw, rounded corners, minimum width=3.0cm, minimum height=0.7cm,
              align=center, font=\small},
  arr/.style={-Stealth, thick}
]
\node[box, fill=yellow!30] (bat) {18650 Li-ion\\3.7 V / 2200 mAh};
\node[box, fill=blue!15,  below left=1.2cm and 1.5cm of bat] (tp)  {TP4056\\Charger};
\node[box, fill=green!20, below right=1.2cm and 0.5cm of bat](mt)  {MT3608\\Boost 5 V};
\node[box, fill=orange!20, below=1.2cm of mt]  (esp) {ESP32\\3.3 V LDO};
\node[box, fill=red!15,   below right=1.2cm and 0.5cm of bat, xshift=4cm] (sim) {SIM800L\\3.7–4.2 V direct};
\node[box, fill=gray!15,  below=1.2cm of esp]  (mpu) {MPU6050 / NEO-6M\\3.3 V};
\draw[arr] (bat) -- (tp);
\draw[arr] (bat) -- (mt);
\draw[arr] (bat) -- (sim);
\draw[arr] (mt)  -- (esp);
\draw[arr] (esp) -- (mpu);
\end{tikzpicture}
\caption{WaveGuard V4 Power Distribution Tree}
\label{fig:power}
\end{figure}

\noindent Key design decisions:
\begin{itemize}
  \item \textbf{SIM800L powered directly from the cell} (3.7--4.2~V) to avoid passing
        the 2~A transmit peak through the MT3608 (rated 2~A maximum).
  \item \textbf{1000~µF electrolytic} bulk capacitor placed within 10~mm of the
        SIM800L VCC pin to absorb the transmit-burst current spike.
  \item \textbf{470~µF electrolytic} on the MT3608 output to stabilise the 5~V rail
        during load transients.
\end{itemize}

\section{I2C Communication and Calibration}

The MPU6050 is initialised with:
\begin{itemize}
  \item Digital Low-Pass Filter (DLPF) = 94~Hz bandwidth (register 0x1A, value 0x02)
  \item Accelerometer Full-Scale = $\pm$2~g (register 0x1C, value 0x00)
  \item Sample rate divider = 19 → 50~Hz (register 0x19, value 0x13)
\end{itemize}

\textbf{Calibration procedure:} The buoy is placed on a flat surface and 500 samples
are averaged to obtain the bias vector $[\bar{a}_x^0, \bar{a}_y^0, \bar{a}_z^0]$.
These offsets are stored in ESP32 NVS (non-volatile storage) and subtracted before the
RMS computation. After calibration, stationary readings are $<$0.005~g RMS.

\textbf{Threshold derivation:} Controlled tests in a calm swimming pool gave
$a_{\rm rms} \approx 0.015$--0.025~g for gentle lapping; vigorous hand-generated
waves produced $>$0.08~g. The thresholds 0.03~g and 0.075~g were set with a safety
margin to avoid false positives while capturing moderate swell.

\section{Software Implementation}

\subsection{Firmware Initialisation}

On boot ($\approx$3~seconds total):
\begin{enumerate}
  \item Serial (115200 baud) for debug output.
  \item I2C on GPIO21/22 at 400~kHz; MPU6050 wakeup and configuration.
  \item UART2 (GPIO16/17, 9600 baud) for NEO-6M GPS.
  \item UART (GPIO26/27, 9600 baud) for SIM800L; AT handshake and SMS mode.
  \item WiFi connection to configured SSID ($\approx$5--10~seconds).
  \item Load calibration offsets from NVS.
\end{enumerate}

\subsection{Main Sampling Loop}

The firmware runs a 1\,000~ms cycle:
\begin{enumerate}
  \item Read 50 MPU6050 samples at 20~ms intervals.
  \item Subtract calibration offsets and compute $a_{\rm rms}$ (Eq.~\ref{eq:rms}).
  \item Read GPS fix (non-blocking; uses last valid fix if unavailable).
  \item Classify state and POST JSON payload to FastAPI backend.
  \item If state = \textsc{high-wave} and SMS cooldown expired (30~s), send SMS.
  \item Update onboard LED indicator (green / yellow / red).
\end{enumerate}

\subsection{FastAPI Backend}

The backend provides the following REST endpoints:

\begin{table}[H]
\centering
\caption{FastAPI REST API Endpoints}
\label{tab:api}
\begin{tabular}{lll}
\toprule
\textbf{Method} & \textbf{Path}            & \textbf{Description} \\
\midrule
GET    & \texttt{/readings}          & Paginated sensor history \\
GET    & \texttt{/readings/latest}   & Most recent reading \\
POST   & \texttt{/readings}          & Ingest new sensor record (firmware) \\
GET    & \texttt{/alerts}            & Recent alert events \\
POST   & \texttt{/auth/login}        & Admin JWT login \\
GET    & \texttt{/recipients}        & List SMS recipients (auth) \\
POST   & \texttt{/recipients}        & Add recipient (auth) \\
DELETE & \texttt{/recipients/\{id\}} & Remove recipient (auth) \\
\bottomrule
\end{tabular}
\end{table}

\subsection{React Dashboard}

The dashboard consists of:
\begin{itemize}
  \item \textbf{Live gauge} showing current $a_{\rm rms}$ with colour-coded alert
        status (green/yellow/red).
  \item \textbf{Time-series chart} (recharts) of the last 60 seconds.
  \item \textbf{Map view} (Leaflet.js) showing buoy GPS position.
  \item \textbf{Alert log} listing the last 10 alert events with timestamps.
  \item \textbf{Admin panel} (JWT-protected) for SMS recipient management.
\end{itemize}

\section{Known Issues and Solutions}

\begin{description}
  \item[SIM800L SMS transmission failure:]
        Early prototypes failed to send SMS intermittently because the cell voltage
        dropped below 3.7~V during the 2~A transmit burst.
        \textit{Solution:} Added the 1\,000~µF bulk capacitor and shortened leads
        to SIM800L; SMS delivery improved to $>$95\%.

  \item[ESP32 random restart during WiFi association:]
        The ESP32 occasionally reset when WiFi scanning was active while MPU6050 I2C
        reads were in progress, causing a stack overflow in interrupt context.
        \textit{Solution:} WiFi initialisation is completed before entering the main
        sensor loop; I2C reads are performed in the main task only.

  \item[Serial monitor baud rate mismatch:]
        Initial debug output was garbled because the IDE monitor was set to 9\,600 while
        the firmware used 115\,200 baud.
        \textit{Solution:} Serial.begin(115200) documented as the required monitor baud.
\end{description}
